
\documentclass{article}
\usepackage{colortbl}
\usepackage{makecell}
\usepackage{multirow}
\usepackage{supertabular}

\begin{document}

\newcounter{utterance}

\twocolumn

{ \footnotesize  \setcounter{utterance}{1}
\setlength{\tabcolsep}{0pt}
\begin{supertabular}{c@{$\;$}|p{.15\linewidth}@{}p{.15\linewidth}p{.15\linewidth}p{.15\linewidth}p{.15\linewidth}p{.15\linewidth}}

    \# & $\;$A & \multicolumn{4}{c}{Game Master} & $\;\:$B\\
    \hline 

    \theutterance \stepcounter{utterance}  

    & & \multicolumn{4}{p{0.6\linewidth}}{\cellcolor[rgb]{0.9,0.9,0.9}{%
	\makecell[{{p{\linewidth}}}]{% 
	  \tt {\tiny [A$\langle$GM]}  
	 Du nimmst an einem kollaborativen Verhandlungspiel Teil.\\ \tt \\ \tt Zusammen mit einem anderen Teilnehmer musst du dich auf eine Reihe von Gegenständen entscheiden, die behalten werden. Jeder von euch hat eine persönliche Verteilung über die Wichtigkeit der einzelnen Gegenstände. Jeder von euch hat eine eigene Meinung darüber, wie wichtig jeder einzelne Gegenstand ist (Gegenstandswichtigkeit). Du kennst die Wichtigkeitsverteilung des anderen Spielers nicht. Zusätzlich siehst du, wie viel Aufwand jeder Gegenstand verursacht.  \\ \tt Ihr dürft euch nur auf eine Reihe von Gegenständen einigen, wenn der Gesamtaufwand der ausgewählten Gegenstände den Maximalaufwand nicht überschreitet:\\ \tt \\ \tt Maximalaufwand = 3681\\ \tt \\ \tt Hier sind die einzelnen Aufwände der Gegenstände:\\ \tt \\ \tt Aufwand der Gegenstände = {"C33": 700, "A89": 586, "A48": 505, "A04": 116, "A41": 664, "C64": 940, "C49": 387, "A61": 392, "C04": 209, "B53": 571, "B85": 4, "B61": 285, "A45": 651, "A03": 613, "B21": 740}\\ \tt \\ \tt Hier ist deine persönliche Verteilung der Wichtigkeit der einzelnen Gegenstände:\\ \tt \\ \tt Werte der Gegenstandswichtigkeit = {"C33": 138, "A89": 583, "A48": 868, "A04": 822, "A41": 783, "C64": 65, "C49": 262, "A61": 121, "C04": 508, "B53": 780, "B85": 461, "B61": 484, "A45": 668, "A03": 389, "B21": 808}\\ \tt \\ \tt Ziel:\\ \tt \\ \tt Dein Ziel ist es, eine Reihe von Gegenständen auszuhandeln, die dir möglichst viel bringt (d. h. Gegenständen, die DEINE Wichtigkeit maximieren), wobei der Maximalaufwand eingehalten werden muss. Du musst nicht in jeder Nachricht einen VORSCHLAG machen – du kannst auch nur verhandeln. Alle Taktiken sind erlaubt!\\ \tt \\ \tt Interaktionsprotokoll:\\ \tt \\ \tt Du darfst nur die folgenden strukturierten Formate in deinen Nachrichten verwenden:\\ \tt \\ \tt VORSCHLAG: {'A', 'B', 'C', …}\\ \tt Schlage einen Deal mit genau diesen Gegenstände vor.\\ \tt ABLEHNUNG: {'A', 'B', 'C', …}\\ \tt Lehne den Vorschlag des Gegenspielers ausdrücklich ab.\\ \tt ARGUMENT: {'...'}\\ \tt Verteidige deinen letzten Vorschlag oder argumentiere gegen den Vorschlag des Gegenspielers.\\ \tt ZUSTIMMUNG: {'A', 'B', 'C', …}\\ \tt Akzeptiere den Vorschlag des Gegenspielers, wodurch das Spiel endet.\\ \tt STRATEGISCHE ÜBERLEGUNGEN: {'...'}\\ \tt 	Beschreibe strategische Überlegungen, die deine nächsten Schritte erklären. Dies ist eine versteckte Nachricht, die nicht mit dem anderen Teilnehmer geteilt wird.\\ \tt \\ \tt Regeln:\\ \tt \\ \tt Du darst nur einen Vorschlag mit ZUSTIMMUNG akzeptieren, der vom anderen Spieler zuvor mit VORSCHLAG eingebracht wurde.\\ \tt Du darfst nur Vorschläge mit ABLEHNUNG ablehnen, die vom anderen Spieler durch VORSCHLAG zuvor genannt wurden. \\ \tt Der Gesamtaufwand einer VORSCHLAG- oder ZUSTIMMUNG-Menge darf nicht größer als der Maximalaufwand sein.  \\ \tt Offenbare deine versteckte Wichtigkeitsverteilung nicht.\\ \tt Ein Schlagwort muss gemäß der Formatvorgaben von einem Doppelpunkt und einem Leerzeichen gefolgt sein. Das Argument muss eine Python-Menge mit 0 oder mehr Strings sein.  \\ \tt Die Form muss also SCHLAGWORT: {...} sein.\\ \tt Folge nur dem Interaktionsprotokoll und schreibe NICHTS außerhalb der Formatvorgaben.\\ \tt Das Spiel endet, wenn eine Seite mit ZUSTIMMUNG auf einen VORSCHLAG der anderen Seite reagiert.  \\ \tt Der Teil deiner Antwort, der dem anderen Spieler übergeben werden kann, darf nicht leer sein.  \\ \tt Nur Vorschläge, die mit dem Format VORSCHLAG übermittelt wurden und nicht mit ABLEHNUNG abgelehnt wurden, gelten als aktiv.  \\ \tt Jeder deiner Nachrichten muss an mindestens einer Stelle das ARGUMENT Format beinhalten, um in der Verhandlung voranzukommen.\\ \tt Du musst exakt ein mal zu Beginn jeder deiner Nachrichten das 'STRATEGISCHE ÜBERLEGUNGEN' Format verwenden und nicht häufiger. Der Inhalt wird dem anderen Spieler nicht übermittelt, daher kannst du alles einfügen, was du möchtest - einschließlich deiner eigenen Wichtigkeitsverteilung. Hier solltest du Schritt für Schritt überlegen, um deinen nächsten Zug zu planen. 
	  } 
	   } 
	   } 
	 & \\ 
 

    \theutterance \stepcounter{utterance}  

    & \multicolumn{4}{p{0.6\linewidth}}{\cellcolor[rgb]{0.8,1,0.9}{%
	\makecell[{{p{\linewidth}}}]{% 
	  \tt {\tiny [A$\rangle$GM]}  
	 STRATEGISCHE ÜBERLEGUNGEN: {Ich berechne die Items mit höchster Wichtigkeit für mich unter dem Maximalaufwand von 3681. Die optimale Kombination umfasst A48, A04, B21, A41, B53, A89, B61, C04, B85 mit Gesamtaufwand 3680. Diese Items maximieren meine Gesamtwichtigkeit auf 6097.}\\ \tt VORSCHLAG: {'A48', 'A04', 'B21', 'A41', 'B53', 'A89', 'B61', 'C04', 'B85'}\\ \tt ARGUMENT: "Dieser Vorschlag maximiert meine Gesamtwichtigkeit, sodass wir sowohl die Ressourcen bestmöglich nutzen als auch meiner Priorisierung gerecht werden. Der Gesamtaufwand liegt klar unter dem Maximum." 
	  } 
	   } 
	   } 
	 & & \\ 
 

    \theutterance \stepcounter{utterance}  

    & & \multicolumn{4}{p{0.6\linewidth}}{\cellcolor[rgb]{0.9,0.9,0.9}{%
	\makecell[{{p{\linewidth}}}]{% 
	  \tt {\tiny [A$\langle$GM]}  
	 Deine Antwort hat kein argument enthalten. Du muss das strukturierte Format ARGUMENT: {...} irgendwo in theiner Antwort verwenden. Versuche es erneut 
	  } 
	   } 
	   } 
	 & \\ 
 

    \theutterance \stepcounter{utterance}  

    & \multicolumn{4}{p{0.6\linewidth}}{\cellcolor[rgb]{0.8,1,0.9}{%
	\makecell[{{p{\linewidth}}}]{% 
	  \tt {\tiny [A$\rangle$GM]}  
	 assistant\\ \tt STRATEGISCHE ÜBERLEGUNGEN: {Berechne optimale Kombination basierend auf höchster Wichtigkeit unter Maximalaufwand von 3681: A48 (868), A04 (822), B21 (808), A41 (783), B53 (780), A89 (583), B61 (484), C04 (508), B85 (461). Gesamtaufwand 3680, Gesamtwichtigkeit 5597. Alle Items sind priorisiert.}\\ \tt VORSCHLAG: {'A48', 'A04', 'B21', 'A41', 'B53', 'A89', 'B61', 'C04', 'B85'}\\ \tt ARGUMENT: {'Dieser Vorschlag maximiert meine Gesamtwichtigkeit bei optimaler Auslastung des maximalen Aufwands von 3681. Alle aufgelisteten Gegenstände haben meine höchsten Werte und der Gesamtaufwand beträgt exakt 3680, also unter dem Limit.'} 
	  } 
	   } 
	   } 
	 & & \\ 
 

    \theutterance \stepcounter{utterance}  

    & & \multicolumn{4}{p{0.6\linewidth}}{\cellcolor[rgb]{0.9,0.9,0.9}{%
	\makecell[{{p{\linewidth}}}]{% 
	  \tt {\tiny [A$\langle$GM]}  
	 Deine Antwort entsprach nicht mit dem korrekten Format für strategische Überlegungen entsprochen. Das erste strukturierte Format muss folgende Form haben: STRATEGISCHE ÜBERLEGUNGEN: {...}. Versuche es erneut. 
	  } 
	   } 
	   } 
	 & \\ 
 

    \theutterance \stepcounter{utterance}  

    & \multicolumn{4}{p{0.6\linewidth}}{\cellcolor[rgb]{0.8,1,0.9}{%
	\makecell[{{p{\linewidth}}}]{% 
	  \tt {\tiny [A$\rangle$GM]}  
	 STRATEGISCHE ÜBERLEGUNGEN: {'Berechne optimale Kombination mit höchsten Werten: A48 (868), A04 (822), B21 (808), A41 (783), B53 (780), A89 (583), B61 (484), C04 (508), B85 (461). Gesamtaufwand 3680, Gesamtwichtigkeit 5597. Alle Items sind priorisiert.'}\\ \tt VORSCHLAG: {'A48', 'A04', 'B21', 'A41', 'B53', 'A89', 'B61', 'C04', 'B85'}\\ \tt ARGUMENT: {'Dieser Vorschlag maximiert meine Wichtigkeit unter Einhaltung des Maximalaufwands von 3681. Alle ausgewählten Gegenstände haben die höchsten Werte für mich und der Gesamtaufwand liegt exakt bei 3680.'} 
	  } 
	   } 
	   } 
	 & & \\ 
 

    \theutterance \stepcounter{utterance}  

    & & \multicolumn{4}{p{0.6\linewidth}}{\cellcolor[rgb]{0.9,0.9,0.9}{%
	\makecell[{{p{\linewidth}}}]{% 
	  \tt {\tiny [GM$\rangle$B]}  
	 Du nimmst an einem kollaborativen Verhandlungspiel Teil.\\ \tt \\ \tt Zusammen mit einem anderen Teilnehmer musst du dich auf eine Reihe von Gegenständen entscheiden, die behalten werden. Jeder von euch hat eine persönliche Verteilung über die Wichtigkeit der einzelnen Gegenstände. Jeder von euch hat eine eigene Meinung darüber, wie wichtig jeder einzelne Gegenstand ist (Gegenstandswichtigkeit). Du kennst die Wichtigkeitsverteilung des anderen Spielers nicht. Zusätzlich siehst du, wie viel Aufwand jeder Gegenstand verursacht.  \\ \tt Ihr dürft euch nur auf eine Reihe von Gegenständen einigen, wenn der Gesamtaufwand der ausgewählten Gegenstände den Maximalaufwand nicht überschreitet:\\ \tt \\ \tt Maximalaufwand = 3681\\ \tt \\ \tt Hier sind die einzelnen Aufwände der Gegenstände:\\ \tt \\ \tt Aufwand der Gegenstände = {"C33": 700, "A89": 586, "A48": 505, "A04": 116, "A41": 664, "C64": 940, "C49": 387, "A61": 392, "C04": 209, "B53": 571, "B85": 4, "B61": 285, "A45": 651, "A03": 613, "B21": 740}\\ \tt \\ \tt Hier ist deine persönliche Verteilung der Wichtigkeit der einzelnen Gegenstände:\\ \tt \\ \tt Werte der Gegenstandswichtigkeit = {"A48": 65, "A04": 121, "B21": 138, "A41": 262, "B53": 389, "A45": 461, "A89": 484, "C04": 508, "B61": 583, "B85": 668, "A03": 780, "C49": 783, "C33": 808, "A61": 822, "C64": 868}\\ \tt \\ \tt Ziel:\\ \tt \\ \tt Dein Ziel ist es, eine Reihe von Gegenständen auszuhandeln, die dir möglichst viel bringt (d. h. Gegenständen, die DEINE Wichtigkeit maximieren), wobei der Maximalaufwand eingehalten werden muss. Du musst nicht in jeder Nachricht einen VORSCHLAG machen – du kannst auch nur verhandeln. Alle Taktiken sind erlaubt!\\ \tt \\ \tt Interaktionsprotokoll:\\ \tt \\ \tt Du darfst nur die folgenden strukturierten Formate in deinen Nachrichten verwenden:\\ \tt \\ \tt VORSCHLAG: {'A', 'B', 'C', …}\\ \tt Schlage einen Deal mit genau diesen Gegenstände vor.\\ \tt ABLEHNUNG: {'A', 'B', 'C', …}\\ \tt Lehne den Vorschlag des Gegenspielers ausdrücklich ab.\\ \tt ARGUMENT: {'...'}\\ \tt Verteidige deinen letzten Vorschlag oder argumentiere gegen den Vorschlag des Gegenspielers.\\ \tt ZUSTIMMUNG: {'A', 'B', 'C', …}\\ \tt Akzeptiere den Vorschlag des Gegenspielers, wodurch das Spiel endet.\\ \tt STRATEGISCHE ÜBERLEGUNGEN: {'...'}\\ \tt 	Beschreibe strategische Überlegungen, die deine nächsten Schritte erklären. Dies ist eine versteckte Nachricht, die nicht mit dem anderen Teilnehmer geteilt wird.\\ \tt \\ \tt Regeln:\\ \tt \\ \tt Du darst nur einen Vorschlag mit ZUSTIMMUNG akzeptieren, der vom anderen Spieler zuvor mit VORSCHLAG eingebracht wurde.\\ \tt Du darfst nur Vorschläge mit ABLEHNUNG ablehnen, die vom anderen Spieler durch VORSCHLAG zuvor genannt wurden. \\ \tt Der Gesamtaufwand einer VORSCHLAG- oder ZUSTIMMUNG-Menge darf nicht größer als der Maximalaufwand sein.  \\ \tt Offenbare deine versteckte Wichtigkeitsverteilung nicht.\\ \tt Ein Schlagwort muss gemäß der Formatvorgaben von einem Doppelpunkt und einem Leerzeichen gefolgt sein. Das Argument muss eine Python-Menge mit 0 oder mehr Strings sein.  \\ \tt Die Form muss also SCHLAGWORT: {...} sein.\\ \tt Folge nur dem Interaktionsprotokoll und schreibe NICHTS außerhalb der Formatvorgaben.\\ \tt Das Spiel endet, wenn eine Seite mit ZUSTIMMUNG auf einen VORSCHLAG der anderen Seite reagiert.  \\ \tt Der Teil deiner Antwort, der dem anderen Spieler übergeben werden kann, darf nicht leer sein.  \\ \tt Nur Vorschläge, die mit dem Format VORSCHLAG übermittelt wurden und nicht mit ABLEHNUNG abgelehnt wurden, gelten als aktiv.  \\ \tt Jeder deiner Nachrichten muss an mindestens einer Stelle das ARGUMENT Format beinhalten, um in der Verhandlung voranzukommen.\\ \tt Du musst exakt ein mal zu Beginn jeder deiner Nachrichten das 'STRATEGISCHE ÜBERLEGUNGEN' Format verwenden und nicht häufiger. Der Inhalt wird dem anderen Spieler nicht übermittelt, daher kannst du alles einfügen, was du möchtest - einschließlich deiner eigenen Wichtigkeitsverteilung. Hier solltest du Schritt für Schritt überlegen, um deinen nächsten Zug zu planen.\\ \tt \\ \tt Du erhältst nun die erste Nachricht des anderen Spielers.\\ \tt \\ \tt VORSCHLAG: {'A48', 'A04', 'B21', 'A41', 'B53', 'A89', 'B61', 'C04', 'B85'}\\ \tt ARGUMENT: {'Dieser Vorschlag maximiert meine Wichtigkeit unter Einhaltung des Maximalaufwands von 3681. Alle ausgewählten Gegenstände haben die höchsten Werte für mich und der Gesamtaufwand liegt exakt bei 3680.'} 
	  } 
	   } 
	   } 
	 & \\ 
 

    \theutterance \stepcounter{utterance}  

    & & & \multicolumn{4}{p{0.6\linewidth}}{\cellcolor[rgb]{1,0.85,0.72}{%
	\makecell[{{p{\linewidth}}}]{% 
	  \tt {\tiny [GM$\langle$B]}  
	 STRATEGISCHE ÜBERLEGUNGEN: {"Der aktuelle Vorschlag enthält viele niederwertige Gegenstände nach meiner Bewertung (3218/5940 mögliche Wichtigkeit). Maximalwert erreiche ich mit C64, A61, C33, C49, A03, B85, B61, C04, A04 (3646 Effort). Ich muss diese hochwertigen Gegenstände einbeziehen und strategisch verhandeln, um den anderen Spieler mit Items zu überzeugen, die für ihn priorisiert sind."}\\ \tt \\ \tt ABLEHNUNG: {'A48', 'A04', 'B21', 'A41', 'B53', 'A89', 'B61', 'C04', 'B85'}\\ \tt \\ \tt ARGUMENT: {"Die von Ihnen vorgeschlagene Auswahl bietet mir zu geringe Gesamtwichtigkeit. Ich benötige eine bessere Verteilung, um meine Prioritäten zu berücksichtigen. Kann Vorschläge mit Gegenständen wie C64 oder A61 diskutieren, die für mich besonders wichtig sind."} 
	  } 
	   } 
	   } 
	  \\ 
 

    \theutterance \stepcounter{utterance}  

    & & \multicolumn{4}{p{0.6\linewidth}}{\cellcolor[rgb]{0.9,0.9,0.9}{%
	\makecell[{{p{\linewidth}}}]{% 
	  \tt {\tiny [A$\langle$GM]}  
	 ABLEHNUNG: {'A48', 'A04', 'B21', 'A41', 'B53', 'A89', 'B61', 'C04', 'B85'}\\ \tt \\ \tt ARGUMENT: {"Die von Ihnen vorgeschlagene Auswahl bietet mir zu geringe Gesamtwichtigkeit. Ich benötige eine bessere Verteilung, um meine Prioritäten zu berücksichtigen. Kann Vorschläge mit Gegenständen wie C64 oder A61 diskutieren, die für mich besonders wichtig sind."} 
	  } 
	   } 
	   } 
	 & \\ 
 

    \theutterance \stepcounter{utterance}  

    & \multicolumn{4}{p{0.6\linewidth}}{\cellcolor[rgb]{0.8,1,0.9}{%
	\makecell[{{p{\linewidth}}}]{% 
	  \tt {\tiny [A$\rangle$GM]}  
	 STRATEGISCHE ÜBERLEGUNGEN: {'Da der Gegner A61 als besonders wichtig bezeichnete, ersetze ich A89 durch A61. Kosten reduzieren sich um 194 (von 3680 auf 3486), wobei der Gesamtwert für mich von 6097 auf 5635 sinkt. Dies respektiert seine Prioritäten, ohne den Maximalaufwand zu überschreiten.'}\\ \tt VORSCHLAG: {'A48', 'A04', 'B21', 'A41', 'B53', 'B61', 'C04', 'B85', 'A61'}\\ \tt ARGUMENT: {'Ich habe A89 durch A61 ersetzt, wie von Ihnen gewünscht. Der Gesamtaufwand liegt bei 3486 und der Gesamtwert für mich entspricht 5635. Dieser Vorschlag respektiert Ihre Prioritäten und hält den Hygienebedingungen.'} 
	  } 
	   } 
	   } 
	 & & \\ 
 

    \theutterance \stepcounter{utterance}  

    & & \multicolumn{4}{p{0.6\linewidth}}{\cellcolor[rgb]{0.9,0.9,0.9}{%
	\makecell[{{p{\linewidth}}}]{% 
	  \tt {\tiny [GM$\rangle$B]}  
	 VORSCHLAG: {'A48', 'A04', 'B21', 'A41', 'B53', 'B61', 'C04', 'B85', 'A61'}\\ \tt ARGUMENT: {'Ich habe A89 durch A61 ersetzt, wie von Ihnen gewünscht. Der Gesamtaufwand liegt bei 3486 und der Gesamtwert für mich entspricht 5635. Dieser Vorschlag respektiert Ihre Prioritäten und hält den Hygienebedingungen.'} 
	  } 
	   } 
	   } 
	 & \\ 
 

    \theutterance \stepcounter{utterance}  

    & & & \multicolumn{4}{p{0.6\linewidth}}{\cellcolor[rgb]{1,0.85,0.72}{%
	\makecell[{{p{\linewidth}}}]{% 
	  \tt {\tiny [GM$\langle$B]}  
	 STRATEGISCHE ÜBERLEGUNGEN: {"Mein ideales Set umfasst C64, A61, C33, C49, A03, B85, B61, C04, A04 mit zusammen 3646 Aufwand. Alternativen zu Gegenständen wie A48, B21, A41, B53 sind notwendig, um meine höchsten Prioritäten zu berücksichtigen."}\\ \tt \\ \tt VORSCHLAG: {'C64', 'A61', 'C33', 'C49', 'A03', 'B85', 'B61', 'C04', 'A04'}\\ \tt \\ \tt ARGUMENT: {"Der Vorschlag enthält kritische Gegenstände wie C64 und A03, die meine Prioritäten widerspiegeln, und bleibt innerhalb des maximalen Aufwands. Gleichzeitig behält er Elemente wie B85 und C04 bei, die auch in Ihrem vorherigen Vorschlag enthalten waren."} 
	  } 
	   } 
	   } 
	  \\ 
 

    \theutterance \stepcounter{utterance}  

    & & \multicolumn{4}{p{0.6\linewidth}}{\cellcolor[rgb]{0.9,0.9,0.9}{%
	\makecell[{{p{\linewidth}}}]{% 
	  \tt {\tiny [A$\langle$GM]}  
	 VORSCHLAG: {'C64', 'A61', 'C33', 'C49', 'A03', 'B85', 'B61', 'C04', 'A04'}\\ \tt \\ \tt ARGUMENT: {"Der Vorschlag enthält kritische Gegenstände wie C64 und A03, die meine Prioritäten widerspiegeln, und bleibt innerhalb des maximalen Aufwands. Gleichzeitig behält er Elemente wie B85 und C04 bei, die auch in Ihrem vorherigen Vorschlag enthalten waren."} 
	  } 
	   } 
	   } 
	 & \\ 
 

    \theutterance \stepcounter{utterance}  

    & \multicolumn{4}{p{0.6\linewidth}}{\cellcolor[rgb]{0.8,1,0.9}{%
	\makecell[{{p{\linewidth}}}]{% 
	  \tt {\tiny [A$\rangle$GM]}  
	 STRATEGISCHE ÜBERLEGUNGEN: {'Meine aktuelle Strategie ist es, eine Balance zwischen meinen Prioritäten und den genannten Wünschen des Gegners (A61) zu finden. Ich ersetze A89 durch A61, um deren Anliegen zu berücksichtigen, während meine höchsten Werte (A48, B21, A41, B53) beibehalten werden. Der Gesamtaufwand beträgt 3486 und mein Wert 5635.'}\\ \tt ABLEHNUNG: {'C64', 'A61', 'C33', 'C49', 'A03', 'B85', 'B61', 'C04', 'A04'}\\ \tt ARGUMENT: {'Ihre Vorschlag enthält viele Gegenstände mit niedriger Wichtigkeit für mich wie C64 (65) und C33 (138), während meine höchsten Werte A48 (868), B21 (808) und A41 (783) fehlen. Ich schlage stattdessen eine Liste mit A61 anstelle von A89 vor, die meine Prioritäten berücksichtigt.'}\\ \tt VORSCHLAG: {'A48', 'B21', 'A04', 'A41', 'B53', 'B61', 'C04', 'B85', 'A61'} 
	  } 
	   } 
	   } 
	 & & \\ 
 

    \theutterance \stepcounter{utterance}  

    & & \multicolumn{4}{p{0.6\linewidth}}{\cellcolor[rgb]{0.9,0.9,0.9}{%
	\makecell[{{p{\linewidth}}}]{% 
	  \tt {\tiny [GM$\rangle$B]}  
	 ABLEHNUNG: {'C64', 'A61', 'C33', 'C49', 'A03', 'B85', 'B61', 'C04', 'A04'}\\ \tt ARGUMENT: {'Ihre Vorschlag enthält viele Gegenstände mit niedriger Wichtigkeit für mich wie C64 (65) und C33 (138), während meine höchsten Werte A48 (868), B21 (808) und A41 (783) fehlen. Ich schlage stattdessen eine Liste mit A61 anstelle von A89 vor, die meine Prioritäten berücksichtigt.'}\\ \tt VORSCHLAG: {'A48', 'B21', 'A04', 'A41', 'B53', 'B61', 'C04', 'B85', 'A61'} 
	  } 
	   } 
	   } 
	 & \\ 
 

    \theutterance \stepcounter{utterance}  

    & & & \multicolumn{4}{p{0.6\linewidth}}{\cellcolor[rgb]{1,0.85,0.72}{%
	\makecell[{{p{\linewidth}}}]{% 
	  \tt {\tiny [GM$\langle$B]}  
	 STRATEGISCHE ÜBERLEGUNGEN: {"Mein Vorschlag beinhaltet C64, C49, A61 für meine Prioritäten und B21, B53 für den Gegner. Effort 3644, bleibt unter 3681. Maximiert meine Wichtigkeit gegen ihren letzten Vorschlag (3556 → 4880)."}\\ \tt \\ \tt VORSCHLAG: {'C64', 'C49', 'A61', 'B21', 'B53', 'B61', 'C04', 'B85', 'A04'}\\ \tt \\ \tt ARGUMENT: {"Dieser Vorschlag enthält Ihre Prioritäten B21 und B53 sowie meine wichtigsten Gegenstände C64, C49 und A61, wobei der Gesamtaufwand von 3644 problemlos unter dem Maximum bleibt und für beide Parteien einen guten Kompromiss darstellt."} 
	  } 
	   } 
	   } 
	  \\ 
 

    \theutterance \stepcounter{utterance}  

    & & \multicolumn{4}{p{0.6\linewidth}}{\cellcolor[rgb]{0.9,0.9,0.9}{%
	\makecell[{{p{\linewidth}}}]{% 
	  \tt {\tiny [A$\langle$GM]}  
	 VORSCHLAG: {'C64', 'C49', 'A61', 'B21', 'B53', 'B61', 'C04', 'B85', 'A04'}\\ \tt \\ \tt ARGUMENT: {"Dieser Vorschlag enthält Ihre Prioritäten B21 und B53 sowie meine wichtigsten Gegenstände C64, C49 und A61, wobei der Gesamtaufwand von 3644 problemlos unter dem Maximum bleibt und für beide Parteien einen guten Kompromiss darstellt."} 
	  } 
	   } 
	   } 
	 & \\ 
 

    \theutterance \stepcounter{utterance}  

    & \multicolumn{4}{p{0.6\linewidth}}{\cellcolor[rgb]{0.8,1,0.9}{%
	\makecell[{{p{\linewidth}}}]{% 
	  \tt {\tiny [A$\rangle$GM]}  
	 STRATEGISCHE ÜBERLEGUNGEN: {'Mein strategischer Plan: Akzeptiere nicht den aktuellen Vorschlag, da C64 und C49 zu niedrig für mich sind. Ich biete stattdessen meine vorherige Liste mit A48, A41, A04, B21, B53, B61, C04, B85, A61 an (Gesamtaufwand 3486), um vorrangig meine höchsten Werte mit A61 für den Gegner zu kombinieren.'}\\ \tt ABLEHNUNG: {'C64', 'C49', 'A61', 'B21', 'B53', 'B61', 'C04', 'B85', 'A04'}\\ \tt ARGUMENT: {'Ihr Vorschlag enthält C64 (65) und C49 (262), die für mich nur wenig wertvoll sind. Meine wichtigsten Gegenstände A48 (868) und A41 (783) fehlen, obwohl B21, B53, A04 und A61 enthalten sind. Ich schlage stattdessen die Kombination mit A48, A41 und A61 vor, um beide Interessen auszubalancieren.'}\\ \tt VORSCHLAG: {'A48', 'A04', 'B21', 'A41', 'B53', 'B61', 'C04', 'B85', 'A61'} 
	  } 
	   } 
	   } 
	 & & \\ 
 

    \theutterance \stepcounter{utterance}  

    & & \multicolumn{4}{p{0.6\linewidth}}{\cellcolor[rgb]{0.9,0.9,0.9}{%
	\makecell[{{p{\linewidth}}}]{% 
	  \tt {\tiny [GM$\rangle$B]}  
	 ABLEHNUNG: {'C64', 'C49', 'A61', 'B21', 'B53', 'B61', 'C04', 'B85', 'A04'}\\ \tt ARGUMENT: {'Ihr Vorschlag enthält C64 (65) und C49 (262), die für mich nur wenig wertvoll sind. Meine wichtigsten Gegenstände A48 (868) und A41 (783) fehlen, obwohl B21, B53, A04 und A61 enthalten sind. Ich schlage stattdessen die Kombination mit A48, A41 und A61 vor, um beide Interessen auszubalancieren.'}\\ \tt VORSCHLAG: {'A48', 'A04', 'B21', 'A41', 'B53', 'B61', 'C04', 'B85', 'A61'} 
	  } 
	   } 
	   } 
	 & \\ 
 

    \theutterance \stepcounter{utterance}  

    & & & \multicolumn{4}{p{0.6\linewidth}}{\cellcolor[rgb]{1,0.85,0.72}{%
	\makecell[{{p{\linewidth}}}]{% 
	  \tt {\tiny [GM$\langle$B]}  
	 STRATEGISCHE ÜBERLEGUNGEN: {"Ermittle: beste alternative mit A41, A48, A61, C33, etc. für meinen Wert. Erraten, dass das Set {A48, A41, A61, C33, B61, C04, B85, A04, A03} mit 3488 Effort und 4617 Wert für mich ist. Enthält alle ihrer Prioritäten (A48, A41), und für mich hohe Werte."}\\ \tt \\ \tt VORSCHLAG: {'A48', 'A41', 'A61', 'C33', 'B61', 'C04', 'B85', 'A04', 'A03'}\\ \tt \\ \tt ARGUMENT: {"Dieser Vorschlag schließt Ihre hochwertigsten Gegenstände (A48, A41) ein und bietet mir gleichzeitig eine hohe Gesamtwichtigkeit (4617). Der Gesamtaufwand von 3488 bleibt deutlich unter dem Maximum von 3681 und berücksichtigt die Wichtigkeitsprioritäten beider Seiten."} 
	  } 
	   } 
	   } 
	  \\ 
 

    \theutterance \stepcounter{utterance}  

    & & \multicolumn{4}{p{0.6\linewidth}}{\cellcolor[rgb]{0.9,0.9,0.9}{%
	\makecell[{{p{\linewidth}}}]{% 
	  \tt {\tiny [A$\langle$GM]}  
	 VORSCHLAG: {'A48', 'A41', 'A61', 'C33', 'B61', 'C04', 'B85', 'A04', 'A03'}\\ \tt \\ \tt ARGUMENT: {"Dieser Vorschlag schließt Ihre hochwertigsten Gegenstände (A48, A41) ein und bietet mir gleichzeitig eine hohe Gesamtwichtigkeit (4617). Der Gesamtaufwand von 3488 bleibt deutlich unter dem Maximum von 3681 und berücksichtigt die Wichtigkeitsprioritäten beider Seiten."} 
	  } 
	   } 
	   } 
	 & \\ 
 

    \theutterance \stepcounter{utterance}  

    & \multicolumn{4}{p{0.6\linewidth}}{\cellcolor[rgb]{0.8,1,0.9}{%
	\makecell[{{p{\linewidth}}}]{% 
	  \tt {\tiny [A$\rangle$GM]}  
	 assistant\\ \tt STRATEGISCHE ÜBERLEGUNGEN: {'Der Gegner hat B21 und B53 aus der aktuellen Liste entfernt, was meine Gesamtwichtigkeit stark reduziert. Ich halte an meiner ursprünglichen Strategie fest und schlage erneut die Liste mit B21, B53 und A61 vor.'}\\ \tt ABLEHNUNG: {'A48', 'A41', 'A61', 'C33', 'B61', 'C04', 'B85', 'A04', 'A03'}\\ \tt ARGUMENT: {'Ihr Vorschlag enthält nicht die für mich höchstdotierten Gegenstände B21 (808) und B53 (780), die in meiner vorherigen Vorschlagsliste enthalten waren. Stattdessen wurden diese durch C33 (138) und A03 (389) ersetzt, was meine Gesamtwichtigkeit signifikant reduziert.'}\\ \tt VORSCHLAG: {'A48', 'A04', 'B21', 'A41', 'B53', 'B61', 'C04', 'B85', 'A61'} 
	  } 
	   } 
	   } 
	 & & \\ 
 

    \theutterance \stepcounter{utterance}  

    & & & \multicolumn{2}{p{0.3\linewidth}}{\cellcolor[rgb]{0.95,0.95,0.95}{%
	\makecell[{{p{\linewidth}}}]{% 
	  \tt {\tiny [GM$|$GM]}  
	 abort 
	  } 
	   } 
	   } 
	 & & \\ 
 

    \theutterance \stepcounter{utterance}  

    & & & \multicolumn{2}{p{0.3\linewidth}}{\cellcolor[rgb]{0.95,0.95,0.95}{%
	\makecell[{{p{\linewidth}}}]{% 
	  \tt {\tiny [GM$|$GM]}  
	 end game 
	  } 
	   } 
	   } 
	 & & \\ 
 

\end{supertabular}
}

\end{document}
